\documentclass[a4paper,11pt,fleqn]{article}%fleqn pour aligner les equations à gauche
\usepackage{../../Styles/Cours}


\newcommand{\chnum}{Chapitre 1}
\newcommand{\chtitle}{Transformation acide-base et pH}
\newcommand{\dtype}{Fiche Révision}

\begin{document}
\maketitle

\section*{Ce qu'il faut connaitre :}
\begin{multicols}{2}
\subsection*{Le vocabulaire :}
    \begin{itemize}[label=$\square$]
    \item Acide et base de Brönsted
    \item Couple acide/base
    \item Espèce amphotère
    \item Réaction acide/base
    \end{itemize}
\subsection*{Les grandeurs :}
    \begin{itemize}[label=$\square$]
        \item pH
    \end{itemize}
\subsection*{Les couples acide/base :}
    \begin{itemize}[label=$\square$]
        \item de l'eau
        \item de l'acide carbonique
        \item des acides carboxyliques
        \item des amines
    \end{itemize}
\subsection*{Les relations :}
    \begin{itemize}[label=$\square$]
        \item entre pH et concentration en ions oxonium
    \end{itemize}
\end{multicols}

\section*{Ce qu'il faut savoir faire :}
\begin{table}[h!]
    \begin{tabular}{|m{.6\linewidth}|m{.2\linewidth}|m{.1\linewidth}|}
        \hline
        \textbf{La compétence} & \textbf{Où la trouver} & \textbf{Je la maîtrise ?} \\
        \hline
        Identifier, à partir d'observations ou de données expérimentales,  un transfert d’ion hydrogène, les couples acide-base mis en jeu et établir l’équation d’une réaction acide-base. & & \\
        \hline
        Représenter le schéma de Lewis et la formule semi-développée d’un acide carboxylique, d’un ion carboxylate, d’une amine et d’un ion ammonium.& & \\
        \hline
        Identifier le caractère amphotère d’une espèce chimique. & & \\
        \hline
        Déterminer, à partir de la valeur de concentration en ion oxonium \ce{H3O+}, la valeur du pH de la solution et inversement.& & \\
        \hline
        Mesurer le pH de solutions d’acide chlorhydrique (\ce{H3O+}, \ce{Cl-}) obtenues par dilutions successives d’un facteur 10 pour tester la relation entre le pH et la concentration en ion oxonium H3O+ apporté.& & \\
        \hline
        Utiliser la fonction logarithme décimal et sa réciproque.& & \\
        \hline
    \end{tabular}
\end{table}

\end{document}
